\section{Introduction}
\label{sec:intro}

Automated root-cause analysis (RCA) for microservice incidents directly affects
time-to-mitigation, and LLM-based agents are the current frontier: they read
heterogeneous telemetry, form hypotheses, drive diagnostic tools, and explain their
conclusions. But an uncomfortable question precedes any accuracy claim for such agents:
\emph{did the agent diagnose the fault, or read the answer off the benchmark?}
Injected-fault datasets --- including ours, initially --- encode their labels
everywhere: run identifiers such as \texttt{tt\_slow\_db\_aggressive\_steady\_r1},
injection containers named \texttt{noisy-neighbor}, metric fields named
\texttt{injection}. A language model that sees any of these needs no telemetry at all.

\paragraph{What we did.} We built (i) \sys, a two-application, four-modality incident
dataset whose kernel-trace layer no existing public dataset provides; (ii) a tool-using
LLM RCA agent operating over deterministic, byte-accounted telemetry tools; and (iii) a
leakage-control harness --- masking, full-fidelity transcripts, and an automated audit
--- that lets us report, to our knowledge for the first time, \emph{certified
hint-free} agent RCA numbers, together with the exact price of the hints we removed.

\paragraph{The evaluation arc is itself a finding.} Our agent initially scored
74/74/61 (service / fault-type / both correct). Closing the leaks collapsed it to
48/17/9 --- proof that it had been leaning on names. Rebuilding the agent with strictly
generic tooling (baseline-vs-incident discipline in every tool, call-graph topology
with peer edges, host and limit-proximity channels, fault-taxonomy definitions)
recovered 83/48/48 --- \emph{above} the leaky number on localization, with every crutch
removed. The sequence quantifies both the contamination risk (26~points of service
accuracy, 57~points of fault classification) and its recoverability by honest means.

\paragraph{Contributions.}
\begin{enumerate}
  \item \textbf{Methodology.} Leakage-controlled agent evaluation: evidence-preserving
  identifier masking (with a cross-tool identity requirement --- masking must not
  fragment one entity into several pseudonyms), ground-truth-free audit transcripts, an
  automated leak auditor, and an A/A variance protocol that bounds single-run claims
  ($\pm$9 points in our setting).
  \item \textbf{Dataset and system.} \sys: 95 labeled runs over Sock~Shop and
  Train-Ticket, four time-aligned modalities (${\sim}0.001$\,ms clock skew), a kernel
  representation ladder (raw CTF $\rightarrow$ per-second KPIs $\rightarrow$ wait
  attribution $\rightarrow$ natural-language digests), pre-registered
  fault$\rightarrow$modality predictions; plus an open agentic-RCA harness with two
  non-LLM baselines and a published state-of-the-art adapter.
  \item \textbf{Empirical results.} (a)~The leak-free agent roughly doubles non-LLM
  baselines on top-1 service localization; (b)~accuracy shows \emph{no cliff} down to
  5\% trace sampling --- traces are ${\sim}20\times$ over-collected for this task;
  (c)~kernel telemetry buys efficiency ($-25$\% agent effort) and induced-DB-latency
  \emph{typing}, while its accuracy-rescue role is absorbed by good cross-modality
  tooling; (d)~\emph{interpreted} kernel summaries beat raw KPIs --- a partial kernel
  view degrades below kernel-blind; (e)~an experience/skill layer is net-negative
  without reliable evidence-based skill selection; unseen-fault incidents under a full
  skill library cost ${\sim}18$ points (a quantified distractor effect, via
  leave-one-family-out).
\end{enumerate}

\todo{one-paragraph roadmap of the paper}
